\section{Why did FAST and Action Binning fail ?}

In the world of LLMs, tokenization is often seen as a backward form of learning
that is redundant. However, the issue is that language is already one of the
most compressed forms of information representation. A single word can pack a
substantial amount of information, so the entire representation is already quite
compressed to begin with.

However, actions are different. A robot operating at 30 Hz can very well be
represented at 15 Hz with some small changes. Then the question becomes:
\emph{Why is compression necessary?}

The idea is that the more compressed the esentation becomes, the easier the
information assignment becomes. In fact, compression leads to shorter generation
horizons, which in turn reduce error accumulation and ultimately lead to better
policies.

\emph{Okay, but FAST already does compression. Why did it fail?} Fair question.
The problem with FAST is its uneven information allocation. Splitting the action
representation into different DCT coefficients leads to the first few tokens
having extremely high entropy, while the suffix tokens essentially encode noise.
This leads to something much worse: due to exposure bias, an early token error
during autoregressive generation can degrade subsequent predictions.
